\documentclass[10pt,a4paper]{amsart}
\usepackage[margin=19mm]{geometry}
\usepackage{amsmath,amssymb,mathtools}
\usepackage[scr=rsfs,cal=euler]{mathalfa}
\usepackage{xfrac}
\usepackage[unicode,hidelinks,bookmarks=false]{hyperref}
\newcommand{\ie}{i.e.\ }
\newcommand{\cf}{cf.\ }
\newcommand{\resp}{respectively\ }
\newcommand{\Subsection}{\subsection}
\providecommand{\nobreakdash}{\nobreak\hbox{-}}
\setlength{\emergencystretch}{2em}
\multlinegap0pt
\usepackage{bm}
\usepackage{bbm}
\usepackage{dsfont}
\usepackage{xspace}
\usepackage{cancel}
\usepackage{mathtools}
\usepackage{amsthm}
\makeatletter
\@ifundefined{thm}{\newtheorem{thm}{Thm}}{}
\@ifundefined{cor}{\newtheorem{cor}[thm]{Corollary}}{}
\@ifundefined{defin}{\newtheorem{defin}[thm]{Definition}}{}
\@ifundefined{lemma}{\newtheorem{lemma}[thm]{Lemma}}{}
\@ifundefined{prop}{\newtheorem{prop}[thm]{Proposition}}{}
\makeatother
\def\zed{{\mathbb Z}}
\newcommand{\Dih}{\operatorname{Dih}}
\newcommand{\bs}{\backslash}
\newcommand{\calG}{\mathcal G}
\newcommand{\calH}{\mathcal H}
\newcommand{\calP}{\mathcal P}
\title[Powell's Conjecture on the Goeritz group of $S^3$ is stably ]{Powell's Conjecture on the Goeritz group of $S^3$ is stably true}
\author{Martin Scharlemann}
\date{Source: arXiv:2210.13629v2 (CC BY 4.0)}
\begin{document}
\maketitle

\noindent\emph{Source and licence.} Powell's Conjecture on the Goeritz group of $S^3$ is stably true. Version arXiv:2210.13629v2 (CC BY 4.0), distributed under the Creative Commons Attribution 4.0 International licence, as recorded in the arXiv licence metadata for this exact version and shown on the abstract page https://arxiv.org/abs/2210.13629v2. Full rights notice: licenses/arxiv-2210.13629-CC-BY-4.0.txt.\par\smallskip

\noindent\textit{Editorial scope.} Counted unit: the Prop whose source label is \texttt{prop:dihedral} in the article (source0.tex lines 962--970, source label \texttt{prop:dihedral}), delivered together with the article's own complete original proof of it (source0.tex lines 972--996, one \texttt{proof} environment of 25 lines). No mathematics is written, repaired or re-derived: statement and proof are the article's own text line for line. Required context retained uncounted: lemma:Kdef (lines 679--687); lemma:homeocommute (lines 734--736); prop:G/P (lines 811--814); defin:thetap (lines 823--826); lemma:fixedp (lines 830--831); lemma:thetaprod (lines 838--839); cor:PpPq (lines 927--929); prop:thetastar (lines 934--940). Every definition, display and label of those lines is retained unchanged. Omitted from this package and not redistributed: 1--678, 688--733, 737--810, 815--822, 827--829, 832--837, 840--926, 930--933, 941--961, 971--971, 997--1196. Nothing inside a retained excerpt is omitted or altered: every retained body is delivered exactly as the article's lines are written, so no omission receipt of any kind accompanies this package. Citations: the retained excerpts carry no citation command at all, so no bibliography entry of the article is transcribed and no reference is redistributed. Reference adaptation: none was needed and none was made; every reference inside the delivered unit resolves inside it, so no unresolved reference or citation is produced. Printed slips kept literal: no printed word, symbol, hypothesis or number of the counted statement or of its proof was altered. Layout and class: the article's own class is \texttt{amsart} with packages including graphicx, amsmath, subfigure, amssymb, color, pinlabel, mathtools, txfonts, hyperref, float, tikz, imakeidx, fancyhdr; the wrapper is the lane's pinned \texttt{amsart} editorial wrapper at 10pt on A4, which restates the theorem-like environments the retained text needs and adds the article's own macro definitions that the retained text needs (\texttt{\textbackslash Dih}, \texttt{\textbackslash bs}, \texttt{\textbackslash calG}, \texttt{\textbackslash calH}, \texttt{\textbackslash calP}, \texttt{\textbackslash zed}), with extra wrapper packages (bm, bbm, dsfont, xspace, cancel, mathtools). The delivered numbering is the unit's own. Assets: no retained span contains a figure environment, a graphics-inclusion command, a PDF-page inclusion, a table or a TikZ picture; no image file is copied into this package, the package declares no asset dependency and no image file is redistributed with it. Licence: version arXiv:2210.13629v2 of this article is distributed under the Creative Commons Attribution 4.0 International licence, as recorded in the arXiv licence metadata for this exact version and shown on the abstract page https://arxiv.org/abs/2210.13629v2. A full rights notice is delivered at licenses/arxiv-2210.13629-CC-BY-4.0.txt. Characters: every retained excerpt is pure ASCII, so the delivered engine, XeLaTeX with its default Latin Modern text font, reports no missing character.
\begin{lemma} \label{lemma:Kdef}  Let $\kappa \subset T_2$ be the $1$-complex that is the union of the 6 boundary circles of the labeled meridians.  
\begin{enumerate}
\item Each of the sets  $\{a_i\}$ and  $\{b_i\}$, $ i \in \zed_3$, of labeled meridians are  setwise invariant under the action of $\Dih_6$ and therefore $\kappa$ is also. 
\item A labeled $A$-meridian is disjoint from exactly one labeled $B$-meridian and meets each of the other two labeled $B$-meridians in a single point (and symmetrically).  
\item A labeled $A$-meridian and a labeled $B$-meridian that meet in a single point are called an {\em orthogonal pair of labeled meridians}; each such pair defines a genus 1 bubble in $T_2$.  
%\item For any orthogonal pair of labeled meridians, say $u_r, w_r$, and any $\psi \in \Dih_6$ the complement in $T_2$ of $(u_r \cup w_r) \cup \psi(u_r \cup w_r)$ is connected.  
\item The complement in the surface $T_2$ of any two orthogonal pairs of labeled meridians is connected. Two examples: $T_2 - [(a_2 \cup b_1) \cup (a_1 \cup b_2)]$ is an annulus and $T_2 - [(a_2 \cup b_1) \cup (a_0 \cup b_2)]$ is a disk. \qed
\end{enumerate} 
\end{lemma}  

\begin{lemma} \label{lemma:homeocommute}  Suppose $p, q \in T_2$, $\psi: (S^3, T_2) \to (S^3, T_2)$ is a homeomorphism, with \[\psi_p: (S^3, T(p)) \to (S^3, T({\psi(p)})) \quad and \quad \psi_q: (S^3, T(q)) \to (S^3, T({\psi(q)}))\] the associated homeomorphism of genus $g$ splittings.  Let $\gamma \subset T_2$ be an arc whose endpoints are $p$ and $q$, so $\psi(\gamma)$ is an arc whose endpoints are  $\psi(p)$ and $\psi(q)$.   Then 
\[ \psi_p \gamma^p_q \sim (\psi(\gamma))^{\psi(p)}_{\psi(q)} \psi_q   :T(q) \to T({\psi(p)}) \]
\end{lemma}

\begin{prop} \label{prop:G/P}  Suppose $p \in T_2 - \kappa$, and $T(p)$ has a semi-standard structure, with $\kappa_f$ the forbidden pair of circles in $\kappa$.  Suppose further that $\psi \in \Dih_6$ and $\gamma_1, \gamma_2$ are arcs in $T_2 - \kappa_f$, each with one end-point on $p$ and the other end point on $\psi^{-1}(p)$.  Then 
\[\psi_{\gamma_1} \calP(p) = \psi_{\gamma_2}\calP(p)\]
%\[\psi_{\psi^{-1}(p)} {\gamma_1}_p^{\psi^{-1}(p)} \calP(p) = \psi_{\psi^{-1}(p)} {\gamma_2}_p^{\psi^{-1}(p)}\calP(p)\]
\end{prop}  

\begin{defin} \label{defin:thetap}  Suppose $p \in T_2 - \kappa$, and $T(p)$ has a semi-standard structure, with $\kappa_f$ the forbidden pair of circles in $\kappa$. 

Let $\Theta_p: \Dih_6 \to \calP(p) \bs\calG(p)/\calP(p)$ be the function given by $\Theta_p(\psi) = \calP(p) \psi_{\psi^{-1}(p)} {\gamma}_p^{\psi^{-1}(p)} \calP(p)$, for $\gamma$ any path in $T_2 - \kappa_f$ that has end points at $p$ and $\psi^{-1}(p)$.  
\end{defin}

\begin{lemma} \label{lemma:fixedp} If $p \in T_2 - \kappa$ is a fixed point for $\psi \in \Dih_6$ then $\Theta_p(\psi) = \calP(p) \psi_p  \calP(p)$.
\end{lemma}

\begin{lemma}   \label{lemma:thetaprod} Suppose, for $p \in T_2 - \kappa$, $\Theta_p (\psi_1) = \calP(p)$ and $\Theta_p (\psi_2) = \calP(p)$.  Then $\Theta_p (\psi_2 \psi_1) = \calP(p)$.
\end{lemma}

\begin{cor} \label{cor:PpPq} For, $p, q \in T_2 - \kappa$, $\Theta_q(\psi) = \alpha_p^q \Theta_p(\psi) \alpha_q^p$ so in particular 
$\Theta_q(\psi) \in \calP(q) \iff \Theta_p(\psi) \in \calP(p)$ and $K(p) = K(q) \subset \Dih_6$.  \qed
\end{cor}

\begin{prop}. \label{prop:thetastar}  Let $\kappa_0 \subset \kappa \subset T_2$ be a pair of circles that bound an orthogonal pair of labeled meridian disks.  One can define, for each $p \in T_2 - \kappa_0$, a collection of homeomorphisms $\calH_p$ from $(S^3, T_g)$ to $(S^3, T(p))$ with the following properties:
\begin{enumerate}
\item For each $p \notin \kappa$, there is a subset $\calH_p^0 \subset \calH_p$ so that each $h^0 \in \calH_p^0$ gives a semi-standard structure on $T(p)$ in which $\kappa_0$ is the forbidden pair of labeled meridians.
\item For each $p \in T^2 - \kappa_0$ and $h_p \in \calH_p$, $\calP_{h_p} = \{\psi \subset \calG(p) | h_p^{-1} \psi h_p \in \calP_g\}$ is independent of $h_p$.  That is, for any pair $h_1, h_2 \in \calH_p$, $\calP_{h_1} = \calP_{h_2}$.  This subgroup of $\calG(p)$ is denoted $\calP(p)$.  
\item For each $p, q \in T_2 - \kappa_0$ and any arc $\gamma \subset T_2 - \kappa_0$, $\calP(q) = \gamma_p^q \calP(p) \gamma_q^p$.
\end{enumerate}
\end{prop}

\begin{prop} \label{prop:dihedral}  Let $\kappa_0 \subset \kappa \subset T_2$ be a pair of circles that bound an orthogonal pair of labeled meridian disks.  One can define, for each $p \in T_2 - \kappa_0$, a function $\Theta_p: \Dih_6 \to \calP(p) \bs \calG(p)/\calP(p)$ with the following properties:
\begin{enumerate}
\item  When $p \notin \kappa$, the definition coincides with that in Definition \ref{defin:thetap}.
\item  For any $p, q \in T_2 - \kappa_0$ and path $\alpha \subset T^2 - \kappa_0$ with end points at $p$ and $q$, $\Theta_q(\psi) = \alpha_p^q \Theta_p(\psi) \alpha_q^p$ so in particular 
$\Theta_q(\psi) = \calP(q) \iff \Theta_p(\psi) = \calP(p)$.
\item  Suppose $\Theta_p (\psi_1) = \calP(p)$ and $\Theta_p (\psi_2) = \calP(p)$.  Then $\Theta_p (\psi_2 \psi_1) = \calP(p)$.
\item Suppose $\psi \in \Dih_6$ has a fixed point $p \notin \kappa_0$.  Then $\Theta_p(\psi) = \calP(p)\psi_p \calP(p)$.   
\end{enumerate}
\end{prop}

\begin{proof} As in the proof of Proposition \ref{prop:thetastar}, pick a fixed base point $\star \in T_2 - \kappa$ and a homeomorphism of pairs $h:(S^3, T_g) \to (S^3, T(\star))$ that gives a semi-standard structure on $T(\star)$ in which $\kappa_0$ is the forbidden pair of labeled meridians.   The function $\Theta_{\star}$ is given in Definition \ref{defin:thetap}.  

For $p \in T^2 - \kappa_0$, choose an arc $\alpha \subset T - \kappa_0$ with end points on $p$ and $\star$ and define $\Theta_p(\psi) =  \alpha_{\star}^p \Theta_{\star}(\psi) \alpha_p^{\star}$.  $\Theta_p(\psi)$ is well-defined, for if $\beta \subset T - \kappa_0$ is another such arc then 
\[ \beta_{\star}^p \theta_{\star}(\psi) \beta_p^{\star} = (\alpha_{\star}^p\alpha^{\star}_p)\beta_{\star}^p \theta_{\star}(\psi) \beta_p^{\star}(\alpha_{\star}^p\alpha^{\star}_p) = \alpha_{\star}^p(\alpha^{\star}_p\beta_{\star}^p) \theta_{\star}(\psi)( \beta_p^{\star}\alpha_{\star}^p)\alpha^{\star}_p = \alpha_{\star}^p \theta_{\star}(\psi)\alpha^{\star}_p\]
The last equality follows because, as argued in the proof of Proposition \ref{prop:G/P}.  the bubble move $\alpha^{\star}_p\beta_{\star}^p \in \calP(\star)$.  

 (1) follows immediately from Corollary \ref{cor:PpPq}.  

The proof of (2) is almost as easy: Suppose $\beta \subset T^2 - \kappa_0$ is a path with end points at $\star$ and $p$, so $\alpha \cup \beta$ has ends at $\star$ and $q$.  Then 
\[ \Theta_q(\psi) = (\alpha \cup \beta)_{\star}^q \Theta_{\star}(\psi) (\alpha \cup \beta)^{\star}_q  =    \alpha_p^q \beta_{\star}^p \Theta_{\star}(\psi) \beta^{\star}_p \alpha_q^p  = \alpha_p^q \Theta_p(\psi) \alpha_q^p \]

Observe that (3) is true for $\Theta_{\star}$ by Lemma \ref{lemma:thetaprod}.  Suppose then that $p \in T_2 - \kappa_0$ and $\Theta_p(\psi_1) = \Theta_p(\psi_2) = \calP(p)$.  Then by (2)
for $i = 1, 2$, $\Theta_{\star}(\psi_i) = \calP(\star)$.  
\[\Theta_p(\psi_2\psi_1) =  \alpha_{\star}^p \Theta_{\star}(\psi_2 \psi_1) \alpha^{\star}_p = \alpha_{\star}^p \calP(\star) \alpha^{\star}_p  = \calP(p),\]
as required.  (The last equality from Proposition \ref{prop:thetastar}(3).)

Claim (4) is true for $p \notin \kappa$ by Lemma \ref{lemma:fixedp}, but our argument will extend to the case $p \in \kappa - \kappa_0$.  Choose $\alpha$ above so that it disjoint not only from $\kappa_0$ but also from $\psi(\kappa_0)$.  This is possible by Lemma \ref{lemma:Kdef}(4).  Since $p$ is a fixed point of $\psi$,  $\gamma = \alpha \cup \psi^{-1}(\alpha)$ is an arc whose endpoints are $\star$ and $\psi^{-1}(\star)$.  Moreover, by choice of $\alpha$, $\gamma$ is disjoint from $\kappa_0$.  Then according to Definition \ref{defin:thetap} 
\[\Theta_{\star}(\psi) = \calP(\star) \psi_{\psi^{-1}(\star)} \gamma_{\star}^{\psi^{-1}(\star)} \calP(\star) = \calP(\star) \psi_{\psi^{-1}(\star)} [\psi^{-1}(\alpha)]_p^{\psi^{-1}(\star)} \alpha_{\star}^p \calP(\star) \]
and so, by  Lemma \ref{lemma:homeocommute}, 
\[\Theta_{\star}(\psi) =  \calP(\star) \alpha_p^{\star} \psi_p \alpha_{\star}^p \calP(\star)  \] 

Hence from Proposition \ref{prop:thetastar}(3)
\[\Theta_p(\psi) = \alpha_{\star}^p \Theta_{\star}(\psi) \alpha_p^{\star} = \alpha_{\star}^p \calP(\star) \alpha_p^{\star} \psi_p \alpha_{\star}^p \calP(\star)  \alpha_p^{\star} = \calP(p)\psi_p \calP(p)\]
as required.
\end{proof}

\end{document}
